\documentclass[english,twoside, openright, hidelinks, 11 pt, class=report,crop=false]{standalone}
\input{../../preamb}
\input{../bm}

\begin{document}

\grubr{opginvsym}
\abc{
\item I \grubrs{opgalgsym2} fant vi at ekstremalverdien $f_\text{e}$ til $f$ er
\[ f_\text{e}=-\frac{b^2}{4a}+c \]
Dermed er
\[ g(x)=f(x)-f_\text{e}=ax^2+bx+c-\left(-\frac{b^2}{4a}+c\right)=ax^2+bx+\frac{b^2}{4a} \]
Ved å løse denne likningen med hensyn på $x$, får vi av \textit{abc}-formelen at
\[ x=\frac{-b\pm\sqrt{b^2-4a\left(\frac{b^2}{4a}-g\right)}}{2a}=\frac{-b\pm2\sqrt{ag}}{2a} \]
Vi lar ${p(g)=x}$ være den omvendte funksjonen til $g$ for $x\in[-\infty, -\frac{b}{2a}]$. Da er
\[ p(g)=\frac{-b-2\sqrt{ag}}{2a} \]
Videre er da
\[ h(g)=p(g)-\left(-\frac{b}{2a}\right)=\sqrt{\frac{g}{a}}\]
\item En endring i $c$ vertikalforskyver grafen, men endrer ikke formen. \os
En endring av $b$ forskyver grafen både horisontalt og vertikalt, men endrer ikke formen. \os
En endring av $a$ forskyver grafen både horisontalt og vertikalt, og endrer formen til grafen. 
}

\newpage
\grubop{opginvsym2}
\abc{
\item Da $A=(a, b)$ ligger på grafen til $f$, er $b=f(a)$. Siden $g$ er den omvendte funksjonen til $f$, vet vi at $g(b)=g(f(a))=a$. Dermed må $(b, a)$ ligge på grafen til $g$, og da er $B=(b, a)$.
\item Vi har at
\[ \vv{AB}=[b-a, a-b] \]
Retningsvektoren $\vec{r}$ til symmetrilinja må stå vinkelrett på $\vv{AB}$. Dette kravet er oppfylt hvis $\vec{r}=[b-a, b-a]$ (se \refops{opgvekrett}). $\vec{r}$ har komponenter med lik lengde og likt fortegn, noe som tilsier at $\vec{r}$ danner en vinkel på $45^\circ$ med $t$-aksen.
}
\fig{opginvsym2_lf}

\newpage
\grubr{opginvareaoftriangle} \\
\abc{
\item Vi setter $B=(b, f(b))$. Da tangenten til $f$ i $B$ danner $60^\circ$ med $x$-aksen, har vi at
\begin{align*}
	f'(b)&=\tan 60^\circ \\
	2b &= \sqrt{3} \\
	b &= \frac{\sqrt{3}}{2}
\end{align*}
Siden $f\left(\frac{\sqrt{3}}{2}\right)=\frac{3}{4}$, er
\[ B=\left(\frac{\sqrt{3}}{2}, \frac{3}{4}\right) \]
Videre er
\[ AD=\frac{DB}{\tan 60^\circ}=\dfrac{3}{4}\cdot\frac{1}{\sqrt{3}}=\frac{\sqrt{3}}{4} \]
Følgelig har vi at
\[ OA=OD-AD=\frac{\sqrt{3}}{2}-\frac{\sqrt{3}}{4}=\frac{\sqrt{3}}{4}  \]
Dermed er
\[ A=\left(\frac{\sqrt{3}}{4}, 0\right) \]
\fig{opginvareaoftriangle_lf_a}
\newpage
\item Av punkt 1) krever vi at
\[ g\left(\frac{\sqrt{3}}{2}\right)=f\left(\frac{\sqrt{3}}{2}\right)=\frac{3}{4} \]
Videre er
{
\small
\begin{align*}
g\left(\frac{\sqrt{3}}{2}\right)&=\frac{(6\sqrt{3} + 8)}{3}\left(\frac{\sqrt{3}}{2}\right)^3 - (4\sqrt{3}+9)\left(\frac{\sqrt{3}}{2}\right)^2 + \frac{(7\sqrt{3} + 8)}{2}\frac{\sqrt{3}}{2} \\
&=\frac{(6\sqrt{3} + 8)}{3}\cdot\frac{\sqrt{3}}{2}\cdot\frac{3}{4} - (4\sqrt{3}+9)\frac{3}{4} + \frac{(7\sqrt{3} + 8)}{2}\frac{\sqrt{3}}{2}\\
&=\frac{(6\sqrt{3} + 8)}{3}\cdot\frac{\sqrt{3}}{2}\cdot\frac{3}{4} - (4\sqrt{3}+9)\frac{3}{4} + \frac{(7\sqrt{3} + 8)}{2}\frac{\sqrt{3}}{2} \\
&=\frac{1}{24}\left(3(18+8\sqrt{3})-18(4\sqrt{3}+9)+6(21+8\sqrt{3})\right) \\
&=\frac{1}{24}\cdot 18 \br
&= \frac{3}{4}
\end{align*}
}

Dermed er kravet oppfylt. \vsk 

Av punkt 2) krever vi at i punktet $B$ er
\[ g''(x)=0 \]
Videre har vi at
\begin{align*}
	g'(x)&=(6\sqrt{3} + 8)x^2 - 2(4\sqrt{3}+9)x + \frac{7\sqrt{3} + 8}{2} \\
	g''(x)&=2(6\sqrt{3} + 8)x - 2(4\sqrt{3}+9)
\end{align*}
I punktet $B$ er
\begin{align*}
	g''\left(\frac{\sqrt{3}}{2}\right)&=2(6\sqrt{3} + 8)\frac{\sqrt{3}}{2} - 2(4\sqrt{3}+9) \\
	&=18+8\sqrt{3}-8\sqrt{3}-18 \\
	&=0
\end{align*}
Følgelig er kravet oppfylt.\vsk

Av vinkelsummen til $\triangle AEB$ har vi at 
\[ \angle BEA=180^\circ-\angle EAB-\angle ABE=180^\circ-60^\circ-45^\circ=75^\circ \]
Av punkt 3) krever vi at i punket $B$ er
\[ g'(x)=-\tan 75^\circ =-2-\sqrt{3} \]

I puntet $B$ er
\begin{align*}
	g'\left(\frac{\sqrt{3}}{2}\right)&=(6\sqrt{3} + 8)\left(\frac{\sqrt{3}}{2}\right)^2 - 2(4\sqrt{3}+9)\frac{\sqrt{3}}{2} + \frac{7\sqrt{3} + 8}{2} \\
	&=(6\sqrt{3}+8)\frac{3}{4}- 2(4\sqrt{3}+9)\frac{\sqrt{3}}{2} + \frac{7\sqrt{3} + 8}{2} \\
	&=\frac{1}{4}\left(18\sqrt{3}+24-(48+36\sqrt{3})+14\sqrt{3}+16\right)\\
	&=\frac{1}{4}(-8-4\sqrt{3}) \\
	&=-2-\sqrt{3}
\end{align*}
Dermed er kravet oppfylt.\newpage


}




\end{document}